% Subsection 4.2.3 - Đặc tả Phương thức
\subsection{Đặc tả Phương thức}
\label{subsec:method-specification}

Phần này trình bày thuật toán xử lý của các phương thức quan trọng nhất. Chi tiết đầy đủ được trình bày trong Phụ lục~\ref{appendix:method-specification}.

\subsubsection{Phương thức createCV}

\begin{enumerate}
    \item Kiểm tra User tồn tại trong hệ thống và số lượng CV hiện tại chưa vượt quá giới hạn 5 CV.
    \item Tạo đối tượng CV mới với các thuộc tính cơ bản: title, fullName, email, phoneNumber, dateOfBirth, gender, address, currentPosition, summary, objective.
    \item Tạo các thành phần con từ dữ liệu đầu vào: Education, WorkExperience, Skill, Project, Certification, Language.
    \item Nếu đây là CV đầu tiên của User, gán \texttt{isMain = true}.
    \item Lưu CV cùng tất cả thành phần con vào cơ sở dữ liệu trong một transaction.
    \item Trả về đối tượng CV vừa tạo.
\end{enumerate}

\subsubsection{Phương thức applyJob}

\begin{enumerate}
    \item Truy vấn Job theo \texttt{jobId}, kiểm tra Job đang ACTIVE và chưa hết hạn (\texttt{expiresAt > now()}).
    \item Kiểm tra CV thuộc về User hiện tại.
    \item Kiểm tra User chưa có Application đang hoạt động (PENDING, REVIEWING, ACCEPTED) cho Job này.
    \item Tạo đối tượng Application với \texttt{status = PENDING}, liên kết với candidate, CV, Job.
    \item Tăng \texttt{applicationCount} của Job thêm 1.
    \item Lưu Application và cập nhật Job vào cơ sở dữ liệu.
    \item Tạo và gửi Notification loại \texttt{APPLICATION\_RECEIVED} đến Recruiter của Company.
    \item Trả về đối tượng Application vừa tạo.
\end{enumerate}

\subsubsection{Phương thức updateApplicationStatus}

\begin{enumerate}
    \item Truy vấn Application theo \texttt{applicationId}.
    \item Kiểm tra người gọi là thành viên của Company sở hữu Job liên quan.
    \item Kiểm tra chuyển trạng thái hợp lệ: PENDING có thể sang REVIEWING hoặc REJECTED; REVIEWING có thể sang ACCEPTED hoặc REJECTED.
    \item Cập nhật \texttt{status} và \texttt{updatedAt} của Application.
    \item Lưu Application vào cơ sở dữ liệu.
    \item Tạo và gửi Notification loại \texttt{APPLICATION\_STATUS\_CHANGED} đến ứng viên.
    \item Trả về đối tượng Application đã cập nhật.
\end{enumerate}

\subsubsection{Phương thức postJob}

\begin{enumerate}
    \item Kiểm tra người gọi là thành viên của Company với vai trò OWNER, MANAGER hoặc RECRUITER.
    \item Kiểm tra Company đang ở trạng thái ACTIVE.
    \item Nếu Job được đăng với trạng thái ACTIVE, kiểm tra \texttt{expiresAt > now()}.
    \item Tạo đối tượng Job với status mặc định DRAFT, \texttt{applicationCount = 0}.
    \item Lưu Job vào cơ sở dữ liệu.
    \item Nếu đăng ACTIVE, tạo và gửi Notification loại \texttt{JOB\_POSTED} đến followers của Company.
    \item Trả về đối tượng Job vừa tạo.
\end{enumerate}

\subsubsection{Phương thức registerCompany}

\begin{enumerate}
    \item Kiểm tra User chưa là thành viên của bất kỳ công ty nào.
    \item Kiểm tra \texttt{documentFile} đã được upload.
    \item Upload documentFile lên Firebase Storage, lấy URL.
    \item Tạo đối tượng Company với \texttt{status = PENDING} và các thông tin từ request.
    \item Tạo đối tượng CompanyMember với \texttt{role = OWNER} liên kết User với Company.
    \item Lưu Company và CompanyMember trong một transaction.
    \item Tạo và gửi Notification loại \texttt{COMPANY\_REGISTRATION} đến tất cả Admin.
    \item Trả về đối tượng Company vừa tạo.
\end{enumerate}

\subsubsection{Phương thức inviteMember}

\begin{enumerate}
    \item Truy vấn Company theo \texttt{companyId}, kiểm tra Company đang ACTIVE.
    \item Kiểm tra người gọi có vai trò OWNER hoặc MANAGER trong Company.
    \item Truy vấn User theo email, kiểm tra User tồn tại và chưa là thành viên của Company.
    \item Tạo đối tượng CompanyMemberInvitation với \texttt{status = PENDING}, \texttt{expiresAt = now() + 7 ngày}.
    \item Lưu Invitation vào cơ sở dữ liệu.
    \item Tạo và gửi Notification loại \texttt{COMPANY\_INVITATION} đến User được mời.
    \item Trả về đối tượng Invitation vừa tạo.
\end{enumerate}
